% ---------------------------------------------------------------------------
% Parte 4 — requer no preâmbulo: \usepackage{booktabs,graphicx}
% Figuras: copiar de code/results/ para figures/:
%   part4_fewshot.png, part4_svm_c.png, part4_examples_dinov2.png
% ---------------------------------------------------------------------------
\section{Representações pré-treinadas: CNN, DINOv2 e CLIP}
\label{sec:part4}

\paragraph{Protocolo.}
Usamos as features congeladas da Sec.~\ref{sec:method}: ResNet-18
\cite{he2016resnet}, ViT-B/32 \cite{dosovitskiy2021vit}, DINOv2 ViT-S/14
\cite{oquab2024dinov2} (CLS e, como extensão, a média dos patch tokens) e
CLIP ViT-B/32 \cite{radford2021clip}. O mapa de features (nenhum, L2,
$\sqrt{\cdot}$+L2) é escolhido no val mAP, e o L2 ganhou em todos os
modelos.

\begin{table*}[t]
\centering
\caption{AP de test por categoria (SVM linear, $C=10$, mapa escolhido no val).
As linhas handcrafted são as escolhas da Parte~2. Extração em CPU por imagem.}
\label{tab:part4}
\small
\resizebox{\linewidth}{!}{%
\begin{tabular}{lrrrrrrrr}
\toprule
Representação & mbike & aero & person & horse & car & mAP & $d$ & Extr.\ (ms) \\
\midrule
BoVW $K{=}512$, Hellinger     & .632 & .691 & .750 & .662 & .734 & .694 &  2048 & 240 \\
VLAD $K{=}32$                 & .707 & .726 & .768 & .735 & .793 & .746 & 10240 & 207 \\
FV $K{=}32$                   & .783 & .759 & .811 & .758 & .807 & .784 & 20608 & 296 \\
\midrule
ResNet-18 layer2              & .767 & .768 & .857 & .725 & .786 & .781 &   128 &  18 \\
ResNet-18 layer3              & .894 & .881 & .963 & .876 & .902 & .903 &   256 &  18 \\
ResNet-18 layer4              & .974 & .968 & .973 & .954 & .954 & .964 &   512 &  18 \\
ViT-B/32 (ImageNet)           & .982 & .987 & .937 & .954 & .945 & .961 &   768 &  35 \\
DINOv2 ViT-S/14, CLS          & .982 & 1.00 & .978 & .994 & .951 & .981 &   384 &  52 \\
DINOv2 ViT-S/14, média patches & .980 & 1.00 & .988 & .985 & .967 & .984 &   384 &  52 \\
CLIP ViT-B/32 + SVM           & .994 & .998 & .991 & .992 & .979 & \textbf{.991} & 512 & 35 \\
CLIP zero-shot (``a photo of a \dots'') & .889 & .982 & .583 & .811 & .865 & .826 & 512 & 35 \\
\bottomrule
\end{tabular}}
\end{table*}

\paragraph{Features aprendidas vs.\ handcrafted (Stage A).}
O salto é grande: a ResNet-18 layer4 obtém 0.964 de mAP, contra 0.784 (FV),
0.746 (VLAD) e 0.694 (BoVW) (Tab.~\ref{tab:part4}). O ganho é maior nas
classes em que o BoVW era pior (motorbike 0.632$\to$0.974, horse
0.662$\to$0.954). No pipeline handcrafted, o descritor local é fixo
(histogramas de gradientes de patches pequenos), o codebook é aprendido sem
labels e o pooling ignora a ordem. Só o SVM vê labels. A CNN aprendeu uma
hierarquia não-linear com 1.28\,M imagens rotuladas do ImageNet, que incluem
classes próximas das nossas, e cada posição do layer4 vê quase a imagem toda
e responde a partes e objetos. Mais dimensão \emph{não} implica melhor AP. O
layer4 (512-d) supera o FV (20\,608-d) por 0.18, o layer3 (256-d) já o supera
com 80$\times$ menos dimensões, e o melhor modelo visual, o DINOv2, tem só
384-d. Dentro da ResNet, o AP sobe com a profundidade (0.781$\to$0.903$\to$0.964),
não por haver mais dimensões \cite{razavian2014,chatfield2014}. A extração
é também mais barata em CPU (18 contra 210--300\,ms por imagem).

\paragraph{DINOv2 (Stage B).}
Com um SVM linear sobre o DINOv2 congelado obtemos 0.981 (CLS) e 0.984 (média
dos patches), acima da ResNet-18 (0.964) e do ViT-B/32 supervisionado
(0.961). A arquitetura é a mesma família, por isso a diferença vem do
pré-treino. O DINOv2 foi treinado sem labels em 142\,M imagens, com objetivos
que obrigam vistas diferentes da mesma imagem a ter o mesmo embedding. Imagens
semanticamente próximas ficam juntas, e categorias como ``cavalo'' formam
grupos quase linearmente separáveis. As invariâncias difíceis (pose,
iluminação, fundo) já estão nas features, e o SVM só tem de encontrar uma
direção, o que pede poucos exemplos (Stage C) \cite{oquab2024dinov2,tian2020}.
Os 256 patch tokens guardam \emph{onde} está cada coisa. Um descritor global
perde a disposição espacial, o número e o tamanho relativo dos objetos e a
separação entre objetos que coexistem (cavaleiro e cavalo). Na média, um
objeto pequeno é diluído pelo fundo. O CLS é um resumo dominado pelo conteúdo
mais saliente. O mAP é semelhante, mas as confusões diferem: com o CLS só 8 de
37 motos com uma única categoria são bem classificadas (18 vão para car),
contra 21 com a média (Tab.~\ref{tab:part4conf}). As imagens de background
classificadas como person descem de 67 para 31. Como os SVMs
one-vs-background não estão calibrados entre si, a comparação é indicativa.

\begin{table*}[t]
\centering
\caption{Matrizes de confusão de test (linhas: verdade; colunas: previsão),
só imagens com uma categoria. DINOv2: SVM com score máximo, ou background se
todos os scores forem negativos. CLIP: prompt mais semelhante entre os 5 da
classe e um de background.}
\label{tab:part4conf}
\small
\setlength{\tabcolsep}{3.5pt}
\begin{tabular}{l rrrrrr c rrrrrr}
\toprule
 & \multicolumn{6}{c}{DINOv2, média patches + SVM} & & \multicolumn{6}{c}{CLIP zero-shot} \\
\cmidrule{2-7}\cmidrule{9-14}
 & mb & ae & pe & ho & ca & bg & & mb & ae & pe & ho & ca & bg \\
\midrule
mbike  & 21 &  0 &   7 &  0 &   6 &    3 & &  36 &   0 &   0 &   1 &   0 &   0 \\
aero   &  0 & 87 &  14 &  0 &  11 &    0 & &   0 & 112 &   0 &   0 &   0 &   0 \\
person &  1 &  1 & 617 &  0 &  15 &   50 & &  26 &  11 & 384 &  14 &  10 & 239 \\
horse  &  0 &  0 &   3 & 24 &   1 &    3 & &   0 &   0 &   0 &  30 &   0 &   1 \\
car    &  1 &  1 &  37 &  0 & 173 &   16 & &   3 &   1 &   0 &   0 & 165 &  59 \\
bg     &  1 &  0 &  31 &  0 &  24 & 1021 & &  31 &  42 & 288 & 106 &  32 & 578 \\
\bottomrule
\end{tabular}
\end{table*}

\begin{figure*}[t]
\centering
\includegraphics[width=\linewidth]{figures/part4_fewshot.png}
\caption{Stage C, DINOv2 CLS + SVM linear: val AP em função do número de
positivos de treino, sem augmentation e com flip horizontal (média $\pm$
desvio em 5 sorteios).}
\label{fig:part4fewshot}
\end{figure*}

\begin{figure}[t]
\centering
\includegraphics[width=\linewidth]{figures/part4_svm_c.png}
\caption{Stage D, DINOv2 CLS + SVM linear: mAP de treino e de validação em
função de $C$, com 5 positivos e com todos.}
\label{fig:part4svmc}
\end{figure}

\paragraph{Few-shot e augmentation (Stage C).}
Usámos o DINOv2 CLS com 1, 5, 10, 20, 50 e todos os positivos, sorteados 5
vezes. Os negativos são sempre todas as imagens de background de treino. A
qualidade da representação reduz a dependência de labels em mais de uma ordem
de grandeza (Fig.~\ref{fig:part4fewshot}). O val mAP é 0.779 com 1 positivo,
0.945 com 5, 0.966 com 20 e 0.980 com todos (90--821). Um único positivo já
supera o BoVW treinado com todos (0.717, Parte~1). A curva satura cedo: de 20
positivos para todos (4--40$\times$ mais labels) ganha-se só $+0.014$.
Continuam a melhorar as classes com mais variação intra-classe, person
(0.949$\to$0.978) e car (0.905$\to$0.945), enquanto o aeroplane chega a 1.000
com 20. Com 1 positivo, o resultado depende da imagem sorteada (motorbike
$\pm0.36$). O flip horizontal prejudica quando o label depende da orientação
(texto, dígitos, sinais, objetos com lado definido). Nas nossas classes
preserva o label e não muda nada: altera o val mAP no máximo 0.005, e no test
(TTA, condição separada) dá 0.9823 sem TTA, 0.9816 com a média dos scores e
0.9818 com a média das features. Plausivelmente porque o DINOv2 já foi
treinado com flips. As duas médias não são equivalentes em geral: para
$s(x)=\mathbf{w}^\top\phi(x)+b$ só coincidem se $\phi$ for linear. Com L2,
para duas features unitárias a um ângulo $\theta$, a média das features
multiplica o termo linear por $1/\cos(\theta/2)$, um fator que varia por
imagem e pode mudar o ranking. Aqui $\theta$ é pequeno.

\paragraph{Escolha de $C$ (Stage D).}
Com 5 positivos, o AP de treino é 1.000 para todo o $C$, porque 5 pontos em
384-d são sempre separáveis. Só a validação revela o sobreajuste
(Fig.~\ref{fig:part4svmc}). O melhor $C$ é o de regularização mais forte
($C=0.1$: 0.955 $\pm$ 0.008). A partir de $C=10$ todas as restrições estão
satisfeitas e a solução não muda (0.945 $\pm$ 0.013). Sem regularização, o
$\mathbf{w}$ é fixado por poucos support vectors e varia muito com a amostra.
Um $C$ pequeno troca variância por bias, e o desvio também desce. Com todos os
positivos aparece a curva clássica: $C=0.1$ sub-ajusta (treino 0.987), $C\ge100$
sobreajusta (treino 1.000, val 0.972), e o ótimo é $C=1$ (0.989). O $C$ ótimo
cresce com os dados. Avaliados uma vez no test, os $C$ escolhidos dão 0.949 (5
positivos) e 0.987 (todos), acima dos 0.981 com o $C=10$ fixo.

\begin{table}[t]
\centering
\caption{CLIP zero-shot: AP de validação por conjunto de prompts. Só o
conjunto escolhido (``photo'') foi avaliado no test.}
\label{tab:part4prompts}
\small
\resizebox{\linewidth}{!}{%
\begin{tabular}{lcrrrrrr}
\toprule
Prompts & n/classe & mbike & aero & person & horse & car & mAP \\
\midrule
nome (``motorbike'')           & 1 & \textbf{.946} & .981 & .596 & \textbf{.949} & .791 & .853 \\
``a photo of a \dots''         & 1 & .916 & 1.00 & \textbf{.656} & .885 & .827 & \textbf{.857} \\
8 templates                    & 8 & .939 & 1.00 & .615 & .892 & \textbf{.834} & .856 \\
sinónimos                      & 2 & .921 & 1.00 & .610 & .884 & .803 & .844 \\
\midrule
test, ``a photo of a \dots''   & 1 & .889 & .982 & .583 & .811 & .865 & .826 \\
\bottomrule
\end{tabular}}
\end{table}

\paragraph{CLIP zero-shot (Stage E).}
O zero-shot é um classificador linear cujo $\mathbf{w}$ é o embedding do texto
do prompt, sem bias e sem ver imagens do dataset. Mede quão bem a imagem
corresponde à frase, não o que a separa do background \emph{deste} dataset. O
SVM usa os mesmos embeddings, mas aprende $\mathbf{w}$ e $b$ com labels, e
adapta-se aos negativos reais e à definição de cada classe (``person'' inclui
multidões e cavaleiros). Por isso chega a 0.991 contra 0.826
(Tab.~\ref{tab:part4}). A diferença é máxima no person (0.583, perto dos
0.477 de um ranking aleatório) e mínima no aeroplane (0.982). Na classificação
a 6 vias (5 classes + um prompt de background, um desvio do enunciado porque
o test contém background), a accuracy é 0.602. Motorbike e horse têm recall
0.97, mas precision de 0.38 e 0.20 (Tab.~\ref{tab:part4conf}). Os falsos
positivos são conceitos vizinhos que não estão nos prompts: bicicletas
$\to$ motorbike, vacas $\to$ horse, pássaros no céu $\to$ aeroplane. Um único
prompt de background também não modela um conjunto aberto (288 background
$\to$ person, 239 person $\to$ background). A formulação do prompt pesa por
classe, mas pouco em média (Tab.~\ref{tab:part4prompts}). O val mAP varia só
entre 0.844 e 0.857, mas de ``nome'' para ``a photo of a\dots'' o person
sobe 0.060 e o horse desce 0.064. No zero-shot o prompt \emph{é} o
classificador, e as semelhanças imagem--texto variam num intervalo estreito.
O ensemble de 8 templates não ganha ao prompt único, ao contrário do
ImageNet \cite{radford2021clip}, plausivelmente porque o erro dominante é o
background, que nenhum template resolve. Escolher os prompts no test daria
uma estimativa otimista ajustada ao ruído, deixaria de ser zero-shot e
tornaria a comparação injusta \cite{cawley2010}. O risco é real: o conjunto
escolhido cai de 0.857 (val) para 0.826 (test), mais do que toda a diferença
entre conjuntos.

\begin{figure*}[t]
\centering
\includegraphics[width=\linewidth]{figures/part4_examples_dinov2.png}
\caption{DINOv2 (média patches) + SVM: verdadeiros positivos com score mais
alto, falsos positivos (background) e falsos negativos, por categoria. O mapa
é a contribuição de cada patch para o score do SVM (parte positiva; média do
mapa + bias = score). Os mapas da ResNet-18 e os exemplos do CLIP estão no
repositório.}
\label{fig:part4examples}
\end{figure*}

\paragraph{Análise visual e casos difíceis (Stage F).}
Os mapas decompõem exatamente o score do SVM pelas posições da feature map
(7$\times$7 na ResNet, um CAM com os pesos do SVM \cite{zhou2016cam};
16$\times$16 no DINOv2). Nos verdadeiros positivos, a ResNet realça sobretudo
o objeto. O DINOv2 faz o mesmo no person e no car, mas no aeroplane e na
motorbike espalha-se pelo céu e pela estrada (Fig.~\ref{fig:part4examples}).
Os falsos positivos mostram dependência de contexto, e as mesmas imagens
enganam os dois modelos: um barco num atrelado na rua (car), um dirigível
(aeroplane). Os mapas dizem com rigor quanto cada posição contribuiu
\emph{para esta decisão} e revelam uso de contexto. Não mostram que a
evidência está no pixel realçado, porque cada posição do layer4 vê quase a
imagem toda e a self-attention mistura os patches. Também não provam
causalidade (seria preciso tapar o objeto ou o fundo e medir o score),
mostram só a parte positiva, com resolução baixa e em poucos exemplos
extremos \cite{adebayo2018}. As classes mais difíceis são car e person (DINOv2
car 0.967; CLIP person 0.583), que têm a maior variação intra-classe e
aparecem muitas vezes como elementos secundários de cenas de rua. As condições
difíceis são objetos pequenos ou distantes, no bordo ou fora do crop central,
silhuetas em contra-luz e cenas cheias \cite{hoiem2012}. A evidência mais
forte é que modelos treinados de formas muito diferentes falham nas mesmas
imagens: um cavaleiro minúsculo numa praia e uma mota que não aparece no crop
falham na ResNet, no DINOv2 e no CLIP. Quando modelos independentes concordam
num erro, a causa está na imagem e não nas features. A co-ocorrência é a
segunda causa: só 37 das 125 motos e 31 dos 148 cavalos de test têm uma única
categoria, e os SVMs one-vs-background nunca viram as outras classes como
negativos, o que explica as confusões motorbike $\to$ car/person
\cite{torralba2011}.
